\chapter{Tonometry (Eye Pressure Test)}

\section*{What Is This Test?}
Tonometry measures intraocular pressure (IOP), the pressure of fluid inside the eye. It is an important part of glaucoma evaluation because persistently elevated pressure can damage the optic nerve, although glaucoma can occur with pressure in the usual range and elevated pressure does not by itself prove glaucoma.

\section*{Alternative Names}
\begin{itemize}
  \item Eye Pressure Test
  \item Intraocular Pressure Measurement
  \item Glaucoma Pressure Test
  \item Applanation Tonometry
  \item Noncontact Tonometry
\end{itemize}
\section*{Why It Is Ordered}
Glaucoma screening or monitoring; follow-up of suspicious optic-nerve or visual-field findings; assessment after eye injury or surgery; and monitoring treatment intended to lower eye pressure.

\section*{How This Test Is Performed}
For contact applanation tonometry, numbing eye drops are used and a small instrument gently measures the force needed to flatten a tiny area of the cornea. Noncontact tonometry uses a brief puff of air and usually does not require numbing drops. The measurement takes seconds for each eye.

\section*{Preparation, Risks, and Limitations}
Remove contact lenses as directed and tell the eye-care professional about eye infection, recent eye surgery, or corneal disease. Contact testing may cause brief irritation; noncontact testing may startle the patient. Corneal thickness, shape, scarring, and measurement technique affect the result. Tonometry must be combined with optic-nerve examination, visual-field testing, and other eye assessments when glaucoma is suspected.

\section*{References}
\begin{enumerate}
  \item American Academy of Ophthalmology. Glaucoma diagnosis and eye-pressure testing.
  \item National Eye Institute. Glaucoma Tests.
\end{enumerate}

\clearpage
\section*{Understanding Results and Follow-up}
The laboratory report should identify the specimen, method, units, reference interval or decision limit, and whether the result is positive, negative, abnormal, or inconclusive. Reference intervals vary with age, sex, pregnancy, specimen type, assay, and laboratory; a result outside a range is not automatically a diagnosis, and a result within a range does not always exclude disease. Discuss unexpected results with the ordering clinician, who can decide whether repeat or confirmatory testing, treatment, or referral is appropriate. Do not change medicines or delay urgent care based on an isolated result.


\section*{Regulatory Status and Authorization Date}
This chapter describes a test category, not a specific commercial product. FDA, CDSCO, EU IVDR, and MHRA approval or registration dates therefore cannot be assigned without the assay manufacturer, product name, instrument, and intended use. For laboratory-developed tests, an individual agency approval date may not exist. See the Preface for the interpretation of regulatory status.

\clearpage
